% ===================================================================
%  TABELLA nr. 16 — Il grond magasin. --- Il basar. Ils giugarels.
%  Traduction 2026 d'apres texto/io, controlee sur texto/fr, texto/ca
%  et texto/oc. Consigne : voir texto/rm/10-tabella-01.tex.
%
%  DERNIER TABLEAU DE LA COLONNE, ET LA MENAGERIE DE CARTON DE
%  L'ALINEA 3 REPOND A LA COLONNE LITUANIENNE, QUI FINISSAIT SUR LE
%  MEME ETAL.
%
%  Le tableau 16 lituanien avait trouve que le chameau,
%  « kupranugaris », est un compose lituanien pour une bete
%  qu'aucun Lituanien n'a vue chez lui, et que c'est L'EVANGILE qui
%  l'avait nommee : le lieu d'apprentissage peut etre un LIVRE, et un
%  livre en particulier.
%
%  Le romanche, sur le meme etal, garde ses betes du pays —
%  « vulp », « luf », « levra », « murinera » — et prend tout le
%  reste tel quel : « elefant », « rinoceros », « ipopotam »,
%  « crocodil », « dromedar », « panthera », « hiena », « struzz »,
%  « papagei », « boa », « chamel ». Pas un compose.
%
%  POURQUOI L'ECART AVEC LE LITUANIEN ? Parce que ces betes-la ne
%  sont pas entrees par un livre, ELLES SONT ENTREES PAR CE MAGASIN.
%  Le tableau le dit lui-meme : une menagerie de carton, des oiseaux
%  mecaniques, un requin de baudruche. C'est la que l'enfant d'une
%  vallee apprend le nom des animaux qu'il ne verra jamais — et le
%  jouet, a la difference du livre, arrive de Nuremberg ou de Paris
%  avec son etiquette deja imprimee.
%
%  LE LITUANIEN AVAIT DONC RAISON SUR LE PRINCIPE ET LE ROMANCHE LE
%  COMPLETE SUR LE CAS : LE LIEU D'APPRENTISSAGE DECIDE DE LA FORME
%  DU MOT, ET UN LIEU D'APPRENTISSAGE PEUT ETRE UN OBJET FABRIQUE.
%  Un texte se traduit — d'ou le compose lituanien ; un jouet ne se
%  traduit pas — d'ou l'emprunt romanche. C'est, une derniere fois,
%  la regle des tableaux 10 et 15 : c'est la route de la chose qui
%  choisit le mot.
%
%  ET C'EST LE FIL DE TOUTE LA COLONNE, QU'ON PEUT MAINTENANT DIRE
%  D'UNE PHRASE : le romanche ne choisit ni par prestige, ni par
%  parente, ni par frontiere ; IL PREND LE MOT PAR OU LA CHOSE EST
%  ARRIVEE. Le velo par le nord (t2), la solive nulle part parce
%  qu'il en avait deja une (t5), « piz » et « lavina » descendus vers
%  la plaine (t9), le gouvernail monte du sud (t10), le chemin de fer
%  nomme a la maison parce que fait a la maison (t12), l'hotel nomme
%  pour ses clients (t13), la confiserie rapportee de l'etranger par
%  ceux qui l'y tenaient (t14), la pomme de terre venue par les cols
%  (t15), et ici l'elephant venu par un rayon de jouets.
%
%  ECART DECLARE DU BLOC t16-15-1, LE TROISIEME ET DERNIER DES TROIS.
%  Le fac-simile ido y ecrit « mi-sferi d) », sans la parenthese
%  ouvrante. Comme au tableau 8, la coquille est du compositeur de
%  l'ido et non du texte romanche : la parenthese est retablie, et
%  renvoji.py passe l'ecart sans rien dire. UNE TRADUCTION N'EST PAS
%  UNE TRANSCRIPTION.
% ===================================================================

%%K t16-apar-1 apar
\VUsaut{4.0mm}
\VUtitre{15.0pt}{60}{TABELLA nr. 16}
\VUsaut{2.0mm}
\VUpk{13.2pt}{40}{Il grond magasin. --- Il basar.}
\VUpk{13.2pt}{40}{Ils giugarels.}
\VUsaut{3.4mm}

%%K t16-01-1 p
1. --- Il Nouv Onn s'avischinescha. Ils gronds magasins han preparà lur étalages da
\VUgras{giugarels} \textsuperscript{(1)} e da regals da Nouv Onn.

%%K t16-01-2 p
Jau hai rugà mes tat da ma manar en il basar per vesair quella bella exposiziun, ed
el ha consentì.

%%K t16-02-1 p
2. --- L'emprim essan nus restads avant bellas panoplias, che cuntegnan in
\VUgras{sclop}, in \VUgras{kepi}, ina \VUgras{sabla}, in \VUgras{ceinturon} ed in pèr
d'\VUgras{epaulettas}, u in \VUgras{elm}, ina \VUgras{curassa} \textsuperscript{(3)} ed ina
\VUgras{pistola} \textsuperscript{(4)}. Tut quai m'ha interessà bler dapli che las
\VUgras{lantiarnas magicas} \textsuperscript{(2)}.

%%K t16-tit-1 sub
\VUsaut{3.0mm}
\VUpk{10.2pt}{40}{Bellas vistas.}
\VUsaut{2.6mm}

%%K t16-03-1 p
3. --- En il medem rayon era ina \VUgras{sentinella} \textsuperscript{(5)} en sia
\VUgras{chaschetta da guardia}; el pareva far la guardia avant ina entira menagaria da
carton. Quantas bes-chas ch'i era là: in \VUgras{papagei} \textsuperscript{(6)} sin sia
staunga, in \VUgras{rinoceros} \textsuperscript{(7)} cun in corn sin il nas, in \VUgras{ren}
\textsuperscript{(8)}, in \VUgras{struzz} \textsuperscript{(9)}, ina \VUgras{schimgia}
\textsuperscript{(10)} che rusgava ina nusch, ina \VUgras{vulp} \textsuperscript{(11)} che
purtava davent ina giaglina, in \VUgras{crocodil} \textsuperscript{(12)} che scharlattava cun
enormas missellas, in \VUgras{chamel} \textsuperscript{(13)}, in elegant \VUgras{levrier}
\textsuperscript{(14)}, ina \VUgras{panthera} \textsuperscript{(15)}, lura duas bes-chas
ch'jau n'enconuscheva betg e ch'èn, sco ch'ins di, ina \VUgras{belletta}
\textsuperscript{(16)} ed ina \VUgras{hiena} \textsuperscript{(17)}. Là era era in
\VUgras{leopard} \textsuperscript{(18)} ed in \VUgras{luf} \textsuperscript{(21)}; fitg dasperas
eran graziusas \VUgras{rattinas} \textsuperscript{(19)} e mintgamai pitschnas \VUgras{mieurinas}
\textsuperscript{(20)} da gumma. A la fin han in \VUgras{ipopotam} \textsuperscript{(22)} ed in
\VUgras{dromedar} \textsuperscript{(23)} cun in bocel finì la collecziun. Jau fiss restà là
anc pliras uras, sch'i na fiss betg stà dasperas in splendid champ plain
\VUgras{schuldads da plumb} \textsuperscript{(29)}, ch'ha tratg mia attenziun.

%%K t16-tit-2 sub
\VUsaut{4.0mm}
\VUpk{10.2pt}{40}{Schuldads e guerra.}
\VUsaut{2.8mm}

%%K t16-04-1 p
4. --- Mes tat, ch'è \VUgras{invalid} \textsuperscript{(24)} e ch'ha stuì bandunar l'armada
pervi d'ina \VUgras{blessura} ch'al ha procurà la \VUgras{medaglia militara}, ha fitg
gugent da ma raquintar las \VUgras{campagnas} vi da las qualas el ha prendì part. El era
cuntent da m'explitgar ils differents muvaments da la pitschna armada en miniatura.
Ina gruppa da vairs schuldads che visitava il basar --- in \VUgras{schuldà dal geni}
\textsuperscript{(25)}, in \VUgras{chatschader alpin} \textsuperscript{(26)} cun in
\VUgras{béret} sin il chau, in \VUgras{sergent-major} \textsuperscript{(27)} da l'infantaria
ed in \VUgras{maréchal des logis} \textsuperscript{(28)} da l'artillaria cun in \VUgras{galun}
d'aur sin la maniya --- è restada sper nus per ascultar las explicaziuns.

%%K t16-05-1 p
5. --- Mes tat, dal tut fier d'avair in public uschè brilliant, ha improvisà in entir
plan da battaglia.

%%K t16-05-2 p
La citad fortifitgada, cun ses \VUgras{mirs} e sias \VUgras{fossas}, era assediada da
l'\VUgras{armada inimica}, che, suenter in lung \VUgras{bumbardament}, era pronta per
l'\VUgras{assaut}. Ma ils \VUgras{assedids}, sforzads da la fam, han empruvà ina
\VUgras{sortida} vers il \VUgras{champ} dals \VUgras{assediants}. Suenter avair sbattì
l'\VUgras{attatga} en in \VUgras{cumbat} obstinà enturn las \VUgras{tendas}, han ils
assediants \VUgras{victorius} lantschà lur \VUgras{escadruns} per persequitar ils
assaltants \VUgras{battids}. Quels han fatg \VUgras{retratga}, cuvrind il \VUgras{chomp da
battaglia} cun \VUgras{morts} e \VUgras{blessads}. \VUgras{Brancardiers} han purtà quels
ultims en l'\VUgras{ambulanza} per als laschar tgirar dal \VUgras{chirurg}.

%%K t16-05-3 p
Malgrà la resistenza eroica da intgins \VUgras{bataglions} ch'avevan furmà in
\VUgras{quadrat}, era la \VUgras{sconfitta} cumpletta, il rest da l'armada ha prendì la
\VUgras{fugia}, laschond blers \VUgras{prischuniers}. L'inimi è ì a cumplenir sia
\VUgras{victoria} occupond la citad. Forsa vegn quai a esser la fin da la campagna e,
suenter in \VUgras{armistizi}, vegn ins a pudair suttascriver in \VUgras{contract da
pasch}.

%%K t16-tit-3 sub
\VUsaut{4.4mm}
\VUpk{10.2pt}{40}{Regals per Nouv Onn.}
\VUsaut{2.8mm}

%%K t16-06-1 p
6. --- Ils giuvens schuldads, ch'han ris ascultond il pittoresc raquint dal veteran,
al han dumandà intginas autras explicaziuns. Quant a mai, sun jau ì enturn la vendita
per guardar ina bella \VUgras{arcbalestra} \textsuperscript{(32)} ed in \VUgras{arch}
\textsuperscript{(33)} cun sia \VUgras{frizza} \textsuperscript{(31)}. Là hai jau chattà ina
autra collecziun da bes-chas: dus \VUgras{utschels chantadurs} \textsuperscript{(34)} --- in
\VUgras{rusigneul} automatic ed in \VUgras{sgrischun} che chantavan a plaina gula --- ed
ina seria da bes-chas curiusas: ina \VUgras{murinera} \textsuperscript{(35)}, in \VUgras{boa}
\textsuperscript{(36)}, ina \VUgras{tartaruga} \textsuperscript{(37)}, in \VUgras{chaun da mar}
\textsuperscript{(38)}, ina \VUgras{balena} \textsuperscript{(39)} ed in \VUgras{squal}
\textsuperscript{(40)} gunflabel.

%%K t16-07-1 p
7. --- Jau na sun betg restà ditg a las contemplar, perquai ch'jau hai vis quai
ch'jau avev'insumnià: in bellissim \VUgras{teater da marionettas} \textsuperscript{(41)} cun
splendidas \VUgras{decoraziuns} ed ina charmanta \VUgras{cortina} cotschna. Sin la
\VUgras{scena} paravan dus acturs giugar ina \VUgras{rolla} en in \VUgras{toc} fitg
interessant, perquai ch'ils pitschens \VUgras{spectaturs} da carton, installads en las
loschas u sesids sin las \VUgras{sutgas} ed en il \VUgras{parterre} davos il
\VUgras{manader d'orchester}, applaudivan vivamain.

%%K t16-07-2 p
Bedauraivlamain custava quel teater fitg char.

%%K t16-08-1 p
8. --- Dal rest era grev da tscherner ils giugarels, perquai ch'en las vitrinas eran
amassads giugarels da tut las sorts: \VUgras{Arlecchin} \textsuperscript{(42)},
\VUgras{Pulcinella} \textsuperscript{(43)}, \VUgras{Pierrot} \textsuperscript{(44)},
\VUgras{tschuettas} \textsuperscript{(45)} che sbattevan las alas, \VUgras{merlots}
\textsuperscript{(46)} che sivlavan, \VUgras{pudels} che blevan, in \VUgras{fonograf}
\textsuperscript{(47)} e \VUgras{gieus da pazienza}. In da quels giugarels m'ha fitg
divertì: el represchentava ina \VUgras{battaglia navala} \textsuperscript{(48)}, en la quala
ina \VUgras{nav} vegniva attatgada da \VUgras{pirats} sper in pajais plantà cun
\VUgras{palmas da cocos}.

%%K t16-noto-1 noto
\VUnotes{143.2mm}{%
\textit{(*) Vesair la tabella nr.~6.}%
}

%%K t16-09-1 p
9. --- Jau pudeva contemplar fitg facilmain tut quellas maraviglias, perquai ch'i era
mo pauca glieud en il magasin, apaina intgins flâneurs, in \VUgras{dragun}
\textsuperscript{(49)} ed in \VUgras{incassader} \textsuperscript{(50)} che preschentava ina
\VUgras{tratta} a la \VUgras{cassa} \textsuperscript{(51)} principala.

%%K t16-10-1 p
10. --- Jau hai dumandà mes tat a tge ch'servan ils cudeschs mess sin assas davos la
\VUgras{cassiera}; el m'ha respondì ch'quai èn \VUgras{cudeschs da contabilitad}.

%%K t16-tit-4 sub
\VUsaut{3.6mm}
\VUpk{10.2pt}{40}{Gapar enturn la buttega.}
\VUsaut{2.8mm}

%%K t16-11-1 p
11. --- Bandunond il rayon dals giugarels, essan nus ids en la suedlaria per dar in
egliada a las \VUgras{suedlas} \textsuperscript{(53)}, a las \VUgras{staffas}
\textsuperscript{(54)}, a las \VUgras{brinas} \textsuperscript{(55)}, a las \VUgras{morsas}
\textsuperscript{(56)} ed a las \VUgras{scurriadas}. Entant ch'il \VUgras{chaprayon}
\textsuperscript{(57)} dava intginas infurmaziuns a mes tat, sun jau ì a guardar
l'étalage dals \VUgras{porzelains} e dals \VUgras{cristals} \textsuperscript{(58)}, nua
ch'èn bellas \VUgras{scadiolas da salata} e delicatas \VUgras{tekannas}
\textsuperscript{(59)}.

%%K t16-12-1 p
12. --- Per ir sin l'emprim plaun essan nus passads davos la vitrina dals
\VUgras{corsets} \textsuperscript{(60)}, sper la bonnetaria, nua ch'eran exponidas fitg
bellas \VUgras{chautschas} \textsuperscript{(63)}. Dasperas laschava ina \VUgras{vendidra}
\textsuperscript{(61)} tscherner ad ina \VUgras{cumpradra} \textsuperscript{(62)} bellas
\VUgras{stoffas} \textsuperscript{(64)} da seida.

%%K t16-12-2 p
Muntond la stgala hai jau remartgà ch'il magasin è ornà cun \VUgras{lantiarnas}
\textsuperscript{(65)} giapunaisas, cun \VUgras{parasols} \textsuperscript{(66)} e cun enormas
\VUgras{palmas} \textsuperscript{(70)}.

%%K t16-13-1 p
13. --- Sin l'emprim plaun n'era betg bler da vesair; mo mobiglia (*), \VUgras{chaschas
fortas} \textsuperscript{(67)}, \VUgras{commodas} \textsuperscript{(68)}, \VUgras{chars
d'uffants} \textsuperscript{(69)}. Ma la vista generala dal magasin era fitg
\VUgras{divertenta}.

%%K t16-14-1 p
14. --- Sper noss peis hai jau vis ils instruments da musica: \VUgras{arpas}
\textsuperscript{(71)}, \VUgras{guitarras} \textsuperscript{(72)}, \VUgras{mandolinas}
\textsuperscript{(73)}, \VUgras{cornets a pistun} \textsuperscript{(74)}, \VUgras{clarinettas}
\textsuperscript{(75)}, violinas cun lur \VUgras{artgin} \textsuperscript{(76)} ed era ina
\VUgras{gronda tromma} \textsuperscript{(77)}. In pau pli lunsch, tar il rayon dal palpiri,
martgadegiava in \VUgras{student} \textsuperscript{(78)} ina chartatscha cun il
\VUgras{vendider} \textsuperscript{(79)}. Sper els hai jau distinguì in bel
\VUgras{abecedari} \textsuperscript{(80)}.

%%K t16-14-2 p
En il mez dal magasin era erigida ina auta \VUgras{planta da Nadal} \textsuperscript{(81)},
dal tut pervista cun chandailas, cun robettas e cun giugarels da tut las sorts:
\VUgras{chavals da lain} \textsuperscript{(82)}, \VUgras{poppas} \textsuperscript{(83)}, etc.

%%K t16-14-3 p
La \VUgras{parfumaria} \textsuperscript{(84)} cun sias \VUgras{auas da dents} e ses
differents \VUgras{parfums} derasava in fitg bun odur che vegniva fin tar nus.

%%K t16-tit-5 sub
\VUsaut{3.4mm}
\VUpk{10.2pt}{40}{Nus cumprain insatge.}
\VUsaut{2.8mm}

%%K t16-15-1 p
15. --- Perquai che mes tat ha cumenzà a chattar il temp memia lung, essan nus ids
giu cun la gronda stgala. El è restà in mument avant ina \VUgras{tabella d'astronomia}
\textsuperscript{(85)} che represchentava, sco ch'el m'ha ditg, \VUgras{cometas}
\textsuperscript{(a)}, \VUgras{eclipsas da glina e da sulegl} \textsuperscript{(b)}, ina rosa
dals vents cun ils \VUgras{quatter punctgs cardinals} \textsuperscript{(c)} \VUgras{(Nord,
Sid, Ost e Vest)}, ina charta da las duas \VUgras{mesas sferas} \textsuperscript{(d)}, ils
\VUgras{planets} \textsuperscript{(e)} ed il \VUgras{sistem solar} \textsuperscript{(f)}. Jau
n'hai chapì nagut da tut quai, e jau era bler pli interessà da la livrea galunada d'in
\VUgras{giuven da liuraziun} \textsuperscript{(86)} che passava, e dals bels
\VUgras{eventails} \textsuperscript{(87)} ch'eran sper mai, sper ils \VUgras{portamunaidas}
\textsuperscript{(88)} ed ils \VUgras{portafegls} \textsuperscript{(89)}.

%%K t16-16-1 p
16. --- Jau hai era admirà sin l'étalage \VUgras{plimas} \textsuperscript{(90)} e
\VUgras{fivlas da tschinta} \textsuperscript{(91)}, las bellas \VUgras{flurs artifizialas}
\textsuperscript{(94)} graziusamain arrangiadas en chanasters. Las \VUgras{vieulas}, las
\VUgras{violettas da mars}, las \VUgras{hiacintas} \textsuperscript{(95)}, las
\VUgras{geranias} \textsuperscript{(96)}, ils \VUgras{gis} \textsuperscript{(97)} e las
\VUgras{capuzinas} \textsuperscript{(98)} eran fitg bain imitadas.

%%K t16-tit-6 sub
\VUsaut{4.4mm}
\VUpk{10.2pt}{40}{Il regal da mes tat.}
\VUsaut{2.8mm}

%%K t16-17-1 p
17. --- Jau hai rugà mes tat da ma cumprar ina da las \VUgras{biischtas da dents}
\textsuperscript{(92)} ch'eran en ina chascha sper \VUgras{guaglias da chavels}
\textsuperscript{(93)}. El ha consentì e jau sun ì sez a pajar mia cumpra a la cassa,
sper la quala ins preparava ina impurtanta \VUgras{spediziun} \textsuperscript{(100)} per in
\VUgras{client} \textsuperscript{(99)}.

%%K t16-18-1 p
18. --- Dal tut nunspetgadamain è naschida ina lieva malquietezza tranter las
persunas restadas sper ils \VUgras{bronzs} \textsuperscript{(101)} artistics. In
\VUgras{lader} \textsuperscript{(102)} era gist vegnì tschiffà sin fatg d'in
\VUgras{inspectur} \textsuperscript{(103)}, entant ch'el empruvava da rubar insatgi. Ins ha
fatg vegnir in polizist per l'\VUgras{arrestar}, e gist cur ch'nus essan sortids, vegniva
il malfatschent manà en la praschun.

\VUsaut{6.0mm}
\VUfilet{22mm}
